\section{Analytical Model} \label{sec:model}

[One-paragraph overview: we extend the heuristic of \citet{raynes2018sign}, which gives mutator fixation probability as the sum of drift and hitchhiking terms under unlimited adaptive potential, to account for a finite number $K$ of available beneficial mutations.
The extension rests on one observation: once mutators reach the fitness peak they stall, still carrying their deleterious load, so any nonmutator lineage that also establishes at the peak displaces them.
Mutators therefore fix only if they (1) escape their load and (2) hold their lead until nonmutators are gone.
State up front that the model has no fitted parameters.]

\subsection{Parameters and Derived Quantities} \label{sec:model:parameters}

[Sentence introducing notation; refer to \cref{sec:poisson,sec:evolution} for the simulation model the quantities below describe.]

\begin{itemize}
  \item $N$, population size;
  \item $n$, deme size, for populations structured as a 2D lattice of demes;
  \item $x_0$, initial mutator frequency ($x_0 = 1/2$ for 50/50 competitions; $x_0 = 1/N$ for a single invading mutator);
  \item $K$, number of available beneficial mutations (adaptive potential);
  \item $s$, fitness advantage conferred by a one-step higher fitness score ($s = 0.1$);
  \item $\lambda_b$, $\lambda_d$, expected beneficial and deleterious mutations per generation for nonmutators;
  \item $M$, fold increase in mutation rates conferred by the mutator allele ($M = 100$);
  \item $\mu$, per-replication probability that the mutator allele arises \textit{de novo} ($\mu = 10^{-6}$).
\end{itemize}

[Sentence: from these we derive the quantities below. Note that selection in the simulations is rank-based within a tournament of two, so a fitter lineage gains advantage $s$ regardless of the size of its fitness score lead; flag that predictions would differ quantitatively under additive or multiplicative fitness.]

\begin{align}
  \delta &= e^{-\lambda_d} - e^{-M \lambda_d}
  && \text{indirect selection against mutators,} \label{eq:delta} \\
  q_0 &= 1 - \frac{1 - e^{-M \lambda_d}}{s}
  && \text{fraction of mutators free of deleterious mutations,} \label{eq:q0} \\
  p(W) &: \; 1 - p = e^{-W p}
  && \text{establishment probability at mean offspring number } W, \label{eq:establishment} \\
  p_n &= p\big((1 + s)\, e^{-\lambda_d}\big), \quad
  p_m = p\big((1 + s)\, e^{-M \lambda_d}\big)
  && \text{for nonmutators and mutators,} \label{eq:pnpm} \\
  u_n &= \lambda_b\, p_n, \quad
  u_m = M \lambda_b\, q_0\, p_m
  && \text{per-capita supply of established beneficial lineages,} \label{eq:supply} \\
  r &= s - \delta
  && \text{growth rate of an established mutator lineage.} \label{eq:growth}
\end{align}

[Brief justification of each line.
\Cref{eq:delta}: mutators' excess rate of producing loaded offspring, approximately $(M - 1) \lambda_d$.
\Cref{eq:q0}: mutation-selection balance under rank-based selection.
\Cref{eq:establishment}: survival probability of a branching process with Poisson offspring \citep{haldane1927mathematical}, approximately $2 (W - 1)$ for $W$ near one.
\Cref{eq:supply}: only load-free mutators found lasting beneficial lineages; cf.\ \citet{johnson1999beneficial,desai2011balance}.
Give numerical values under simulation parameters: $\delta \approx 0.0099$, $q_0 \approx 0.90$, $p_n \approx 0.176$, $p_m \approx 0.159$, $u_n / u_m \approx 1/81$.]

\subsection{Escape} \label{sec:model:escape}

[Sentence: following \citet{raynes2018sign}, mutators escape their load either by drift or by hitchhiking; we generalize both terms from a single invader to arbitrary start frequency $x_0$.]

\begin{align}
  P_\mathrm{drift} &= \frac{e^{2 N \delta x_0} - 1}{e^{2 N \delta} - 1}, \label{eq:drift} \\
  P_\mathrm{hitch} &= 1 - (1 - x_0)^{N / N_\downarrow},
  \qquad N_\downarrow = \frac{\delta}{u_m}, \label{eq:hitch} \\
  P_\mathrm{escape} &= P_\mathrm{drift} + (1 - P_\mathrm{drift})\, P_\mathrm{hitch}. \label{eq:escape}
\end{align}

[Explain: \cref{eq:drift} is the fixation probability of a deleterious allele under the diffusion approximation \citep{kimura1962probability}.
\Cref{eq:hitch} is the probability that at least one mutator founds an established beneficial lineage while loaded mutators decline at rate $\delta$; for $x_0 = 1/N$ it reduces to eq.\ 1c of \citet{raynes2018sign}.
$P_\mathrm{hitch} = x_0$ exactly at $N = N_\downarrow$, whatever the value of $x_0$: this is the lower critical population size, above which mutators are favored (the sign inversion of \citet{raynes2018sign}); $N_\downarrow \approx 690$ under simulation parameters.
With unlimited adaptive potential, $P_\mathrm{fix} = P_\mathrm{escape}$.]

\subsection{Sweep Time} \label{sec:model:sweep}

[Sentence: the time $T$ for a fitness class to sweep, when lineages of that class are founded at per-capita rate $u$ and grow at rate $r$, is where population structure enters the model, and the only place it does.
Lineages are founded as a Poisson process in space and time and a location stays unswept so long as no earlier-founded lineage has reached it: Kolmogorov--Johnson--Mehl--Avrami kinetics \citep{kolmogorov1937statistical,avrami1939kinetics}, applied to spatial sweeps by \citet{ralph2010parallel,martens2011interfering}.]

\begin{align}
  T_\mathrm{mixed}(u) &= \frac{\ln(p_m\, r / u) - \gamma}{r}, \label{eq:sweep-mixed} \\
  T_\mathrm{2D}(u) &= \Gamma\!\left(\tfrac{4}{3}\right) \left(\frac{3}{\pi\, u\, n\, v^2}\right)^{1/3} + \frac{\ln(n\, p_m)}{r},
  \qquad v = \frac{r}{\ln(n\, r)}. \label{eq:sweep-2d}
\end{align}

[Explain: $\gamma$ is the Euler--Mascheroni constant and $\Gamma$ the gamma function.
In a well-mixed population an established lineage grows exponentially, giving the logarithm in \cref{eq:sweep-mixed}.
On a 2D lattice of demes exchanging one migrant per neighbor per generation, a lineage first takes over its founding deme, then spreads as a disc at wave speed $v$ (demes per generation) \citep{fisher1937wave}; \cref{eq:sweep-2d} assumes strong selection within demes, $n r \gg 1$.
Note that establishment probability itself is unaffected by this form of population structure \citep{maruyama1970fixation}, so structure does not enter $P_\mathrm{escape}$.
In $d$ dimensions the first term of \cref{eq:sweep-2d} scales as $(u v^d)^{-1/(d+1)}$.]

\subsection{Hold} \label{sec:model:hold}

[Sentence: share law. Equal-fitness mutators and nonmutators compete on equal terms, so each new fitness class is divided between them in proportion to the rates at which they found it; cf.\ the theory of recurrent-mutation soft sweeps \citep{pennings2006soft}.]

\begin{align}
  o_k &= \frac{1 - x_0}{x_0}\, P_\mathrm{hitch} \left(\frac{u_n}{u_m}\right)^{k},
  \qquad a_k = \frac{o_k}{1 + o_k}, \label{eq:share} \\
  P_\mathrm{hold} &= \exp\!\left(-N\, u_n\, a_{K-1}\, T_K\right),
  \qquad
  T_K = T\!\left((1 - a_{K-1})\, u_m\right), \label{eq:hold} \\
  P_\mathrm{fix} &= P_\mathrm{escape}\, P_\mathrm{hold}. \label{eq:fix}
\end{align}

[Explain: $o_k$ and $a_k$ are nonmutators' odds and share in fitness class $k = 0, 1, \ldots$; each further class cuts nonmutators' odds by the supply ratio $u_n / u_m$, and $(1 - a_k)\, u_m$ is mutators' supply into class $k + 1$.
The factor $P_\mathrm{hitch}$ conditions on mutators having escaped; it is one for a start at appreciable frequency in a large population, where $a_0 = 1 - x_0$, and matters for a single invader.
Nonmutators one step below the peak can reach it until mutators' final sweep arrives, a time $T_K$ from \cref{eq:sweep-mixed} or \cref{eq:sweep-2d}; the expected number of nonmutator lineages established at the peak in that time is the exponent of \cref{eq:hold}.
This formalizes the verbal argument of \citet{desai2011balance} and the multi-step dynamics shown by \citet{tanaka2003evolution,tenaillon1999mutators}.
Normalized fixation probability, as plotted by \citet{raynes2018sign}, is $P_\mathrm{fix} / x_0$.]

\begin{figure}[h]
\begin{minipage}{0.5\textwidth}
  \includegraphics[width=\linewidth]{bindle/bindle-theory-invader-npfix.py/bindle/teeplots/theory-invader-npfix/viz=plot-invader-npfix+ext=.pdf}%
\end{minipage}%
\begin{minipage}{0.5\textwidth}
  \caption{
  TODO
  }
  \label{fig:invader-npfix}
\end{minipage}
\end{figure}


[Sentence: upper critical population size. For a start at appreciable frequency, $P_\mathrm{hold}$ falls through $x_0$ at]

\begin{equation}
  N_\uparrow(K) = \frac{\ln(1 / x_0)}{u_n\, a_{K-1}\, T_K}, \label{eq:ncrit-upper}
\end{equation}

[Explain: mutators are favored only for $N_\downarrow < N < N_\uparrow(K)$.
Each further available beneficial mutation raises $N_\uparrow$ by the factor $u_m / u_n$, roughly the mutator strength $M$.
With unlimited adaptive potential there is no upper threshold, recovering \citet{raynes2018sign}.
Refer to Figure \ref{fig:invader-npfix} for the single-invader case.
Give predicted values under simulation parameters for a 50/50 start, well-mixed: $N_\uparrow \approx 1.0 \times 10^5$, $4.6 \times 10^6$, $3.7 \times 10^8$, and $3.0 \times 10^{10}$ for $K = 1$, 2, 3, and 4; and for 256-agent demes: $6.7 \times 10^4$, $3.1 \times 10^6$, $2.6 \times 10^8$, and $2.1 \times 10^{10}$.]

\subsection{Mutators Arising \textit{De Novo}} \label{sec:model:denovo}

[Sentence: mutators arising \textit{de novo} are treated as present at their typical frequency in an adapting nonmutator population, $\bar{x}$: their mutation-selection balance frequency $\mu / \delta$ \citep{johnson1999approach}, discounted for being purged by each nonmutator sweep and rebuilding over the interval $T_n$ between sweeps.]

\begin{equation}
  \bar{x} = \frac{\mu}{\delta} \left(1 - \frac{1 - e^{-\delta T_n}}{\delta T_n}\right). \label{eq:xbar}
\end{equation}

[Sentence: small populations, where sweeps do not overlap. Each mutator hitchhikes with probability $u_m / \delta$ before its load removes it, so mutators found sweeps at per-capita rate $\mu\, u_m / \delta$, against $u_n$ for nonmutators.
The probability that a given sweep is a mutator's, and from it fixation probability, are independent of population size:]

\begin{equation}
  \rho = \frac{\mu\, u_m}{\delta\, u_n + \mu\, u_m},
  \qquad
  P_\mathrm{fix} = 1 - (1 - \rho)^K. \label{eq:denovo-small}
\end{equation}

[Sentence: large populations. Fixation probability switches from near zero to near one over a narrow range of adaptive potential; $K^*$ denotes the number of available beneficial mutations at which mutators fix half the time.
Well-mixed: this is \cref{eq:share,eq:hold} at start frequency $\bar{x}$; mutators' odds multiply by $u_m / u_n$ with each class until nonmutators' share below the peak is too small to catch up, $P_\mathrm{hold} = 1/2$.
2D lattice: a mutator lineage that gets ahead adapts faster than its surroundings but spreads no faster, so mutators fix only if their domains, founded at rate $J = n\, \bar{x}\, u_m$ per deme per generation, cover all $N / n$ demes before nonmutators climb $K$ steps at one sweep time $T_n$ each.]

\begin{align}
  K^*_\mathrm{mixed} &= 1 + \frac{\ln\!\left[N\, u_n\, T_\mathrm{mixed}(u_m) / (\bar{x} \ln 2)\right]}{\ln(u_m / u_n)}, \label{eq:kstar-mixed} \\
  K^*_\mathrm{2D} &= \frac{1}{T_n} \left\{\frac{3 \ln\!\left[N / (n \ln 2)\right]}{\pi\, J\, v^2}\right\}^{1/3}, \label{eq:kstar-2d} \\
  \text{peak prevalence} &= 1 - \exp\!\left[-\frac{\pi}{3}\, J\, v^2\, (K\, T_n)^3\right]. \label{eq:peak-prevalence}
\end{align}

[Explain: $T_n$ is \cref{eq:sweep-mixed} or \cref{eq:sweep-2d} evaluated for nonmutators (supply $u_n$, growth $s$, establishment $p_n$).
$K^*$ depends only logarithmically on $N$ and $\mu$, and is far larger under spatial structure.
\Cref{eq:peak-prevalence} gives the fraction of the lattice reached by mutator domains when nonmutators reach the fitness peak.
Note regimes of validity: \cref{eq:kstar-mixed,eq:kstar-2d} assume overlapping sweeps (large $N$); \cref{eq:denovo-small} assumes sequential sweeps (small $N$).]

[Closing paragraph on scope and limitations: closed forms are approximations intended as scaling laws; numerical refinements (diffusion with killing for escape; branching-process recursion for hold; nucleation-and-growth with renewal sweeps for 2D \textit{de novo}) are provided in the supplement/software.
The single-invader predictions of Figure \ref{fig:invader-npfix} have not been checked against simulation.
Selection is rank-based here; a resimulation with additive fitness raised $N_\uparrow$ 6-fold at $K = 2$ and 15-fold at $K = 3$.
Possible relation to clonal interference theory \citep{gerrish1998fate,desai2007beneficial,good2016evolution}.]
